\emph{Context:} AI coding agents now open, trigger and appear in GitHub issues, yet empirical studies of these agents examine mainly their pull requests. \emph{Objective:} This study compares the issues attributed to different agents in volume, state and labelling, time-to-close, and topics. \emph{Method:} A Goal--Question--Metric mining study of \NIssues{} issues from \NRepos{} repositories attributed to \NAgents{} agents, using rank-based and permutation tests with effect sizes, survival analysis that counts open issues, and Ensemble-LDA topics checked for stability and confirmed with BERTopic; attributions, topic labels and the language flag were validated by hand on blind samples. \emph{Results:} Agents differ most in \emph{how} an issue is linked to them (Cram\'er's V = \VAttribution). The agent's own bot wrote \OwnBotShareCoderabbit\% of CodeRabbit and \OwnBotShareDevin\% of Devin issues, and these are closed 18--24 percentage points less often than the same agents' other issues. Most labelled Sweep issues carry only the label that invokes the agent. Time-to-close differs with a medium effect (\(\varepsilon^2 = \KWeps\)), but rankings based on closed issues disagree with survival analysis (Spearman \(\rho = \RankRho\)). Topics depend on the agent under both models, most clearly for CodeRabbit's review follow-ups, though only two of five fine topics were both reproducible and correctly labelled. By hand, \ValCorrectPct\% of attributions were correct, and in half of these the agent was only mentioned. \emph{Conclusions:} The label ``AI agent'' covers distinct kinds of work; agent comparisons should separate the agent's role, measure speed with survival analysis, and validate attributions and topics by hand.
